% !TeX root = ../main.tex
\section{Các nguồn doanh thu dự kiến}

\textbf{Nguồn doanh thu cơ sở (từ Giao dịch B2C):}
\begin{itemize}
    \item \textbf{Phí nền tảng (Take Rate):} Đây là nguồn thu chủ lực của Mutux. Hệ thống sẽ tự động trích xuất hoa hồng \textbf{30\%} trên tổng giá trị tiền thuê của mỗi giao dịch thành công.
\end{itemize}

\textbf{Nguồn doanh thu Dịch vụ (Theo lộ trình 2 giai đoạn):}
\begin{itemize}
    \item \textbf{Dịch vụ Hỗ trợ tiền cọc (Deposit Advance Service):} Trong giai đoạn đầu, Mutux triển khai cơ chế hỗ trợ ứng trước tiền cọc cho một số khách hàng đáp ứng điều kiện xác minh danh tính (eKYC) và tiêu chí đánh giá nội bộ. Nền tảng thu \textbf{phí dịch vụ là 30.000 VNĐ} cho mỗi tháng người dùng sử dụng tiền cọc để nạp vào ví. Dịch vụ này sử dụng mô hình hợp tác với các đối tác tài chính (BNPL như MoMo, Fundiin) để giảm áp lực vốn lưu động và mở rộng khả năng phục vụ. Mutux không trực tiếp cung cấp dịch vụ tín dụng. Trong giai đoạn đầu, nền tảng chỉ đóng vai trò kết nối người dùng với đối tác BNPL hoặc các tổ chức tài chính. Mutux thu hoa hồng giới thiệu hoặc phí dịch vụ từ đối tác. Đối với rủi ro hỏng hóc/mất mát thiết bị, Mutux áp dụng cơ chế Giữ tiền Escrow + Định danh eKYC + Trừ trực tiếp tiền cọc/ví đối soát nếu xảy ra tranh chấp.
\end{itemize}

\textbf{Nguồn doanh thu mở rộng (từ Đối tác B2B - Triển khai sau 1 năm):}
\begin{itemize}
    \item \textbf{Phí quảng cáo và Đẩy tin (Featured Listings):} Nhằm tối ưu hóa khả năng bán hàng cho các shop gaming gear, Mutux thu phí để hiển thị ưu tiên sản phẩm của họ trên trang chủ hoặc khi người dùng tìm kiếm. 
    \item \textbf{Hoa hồng bán chéo (Cross-selling):} Đóng vai trò như một kênh "Dùng thử trước khi mua", Mutux sẽ nhận hoa hồng môi giới khi người thuê quyết định "mua đứt" thiết bị sau quá trình trải nghiệm.
\end{itemize}
